\chapter{Einleitung}
Es existieren viele relevante Probleme mit polynomialer Laufzeit, welche auf vielen Daten angewandt werden müssen für allgemeine Zwecke.
Diese Vorlesung behandelt verschiedene Techniken, welche verwendet werden können, um möglichst viele Daten möglichst schnell zu verarbeiten.
Dabei geht es hauptsächlich nicht um algorithmische Optimierungsschritte, sondern um die vollständige Ausnutzung der Computer Hardware.

Allgemeine Techniken, die in dieser Vorlesung vorgestellt werden sind: Cache- und Speicher-Optimierung, Vektorisierung, Parallelisierung, Verteiltes Rechnen und Rechnen auf Grafikkarten.
